\subsection{Cantor-Lebesgue Singular Function Construction}

\begin{definition} \label{def:cantor-lebesgue-singular-function-source} \hlt{Cantor--Lebesgue Singular Function}\\
Let $P$ be the Cantor set. For $k\geq1$, let $I_{k,1},\ldots,I_{k,2^k-1}$ list, left to right, the components of the stage-$k$ removed set $U_k$. Set $U\defeq\bigcup_{k\geq1}U_k=[0,1]\setminus P$.
\begin{enumerate}[label=(\roman*)]
\setlength{\itemsep}{0pt}
\item For $x\in U$, choose $k\geq1$ and $1\leq j\leq2^k-1$ with $x\in I_{k,j}$, and set $\varphi(x)\defeq\frac{j}{2^k}$.
\item This value is independent of the chosen stage: if $I_{k,j}=I_{\ell,m}$ with $\ell\geq k$, then $m=2^{\ell-k}j$ and $\frac{m}{2^\ell}=\frac{j}{2^k}$.
\item On $P$, set $\varphi(0)\defeq0$ and $\varphi(x)\defeq\sup\{\varphi(t):t\in U\cap[0,x]\}$ for $x\in P\setminus\{0\}$.
\end{enumerate}
The resulting $\varphi:[0,1]\to[0,1]$ is the \hlt{Cantor--Lebesgue singular function}.
\end{definition}

\begin{figure}[H]
  \centering
  \begin{tikzpicture}[
  x=7.6cm,
  y=3.25cm,
  >=Stealth,
  every node/.style={font=\scriptsize},
  axis/.style={->, black!45, line width=0.45pt},
  tick/.style={black!32, line width=0.35pt},
  plateau/.style={blue!70!black, line width=0.95pt, line cap=round},
  guide/.style={black!14, dashed, line width=0.32pt}
]
  \draw[axis] (-0.02,0) -- (1.04,0) node[right] {$x$};
  \draw[axis] (0,-0.02) -- (0,1.05) node[above] {$\varphi(x)$};

  \foreach \y/\lab in {
    0.125/{1/8},
    0.250/{1/4},
    0.375/{3/8},
    0.500/{1/2},
    0.625/{5/8},
    0.750/{3/4},
    0.875/{7/8},
    1.000/{1}
  } {
    \draw[tick] (-0.008,\y) -- (0.008,\y);
    \node[left=2pt] at (0,\y) {$\lab$};
  }

  \foreach \x/\lab in {
    0/{0},
    0.333333/{1/3},
    0.666667/{2/3},
    1/{1}
  } {
    \draw[tick] (\x,-0.008) -- (\x,0.008);
    \node[below=2pt] at (\x,0) {$\lab$};
  }

  \foreach \a/\b/\h in {
    0.037037/0.074074/0.125,
    0.111111/0.222222/0.250,
    0.259259/0.296296/0.375,
    0.333333/0.666667/0.500,
    0.703704/0.740741/0.625,
    0.777778/0.888889/0.750,
    0.925926/0.962963/0.875
  } {
    \draw[guide] (\a,0) -- (\a,\h);
    \draw[guide] (\b,0) -- (\b,\h);
    \draw[plateau] (\a,\h) -- (\b,\h);
  }

\end{tikzpicture}

  \caption{Horizontal segments show the Cantor--Lebesgue plateaus on the intervals removed through stage $3$.}
  \label{fig:cantor-lebesgue-o3-plateaus}
\end{figure}

\begin{proposition} \label{prop:cantor-singular-function-properties} \hlt{Properties of the Singular Function}\\
Let $P$ be the Cantor set, $U\defeq[0,1]\setminus P$, and $\varphi:[0,1]\to[0,1]$ the Cantor--Lebesgue singular function. Then:
\begin{enumerate}[label=(\roman*)]
\setlength{\itemsep}{0pt}
\item $\varphi$ is nondecreasing.
\item $\varphi$ is continuous on $[0,1]$.
\item $\varphi([0,1])=[0,1]$.
\item $\varphi'(x)=0$ for every $x\in U$.
\item $m(U)=1$.
\end{enumerate}
\begin{proof}
\begin{enumerate}[label=(\roman*)]
\setlength{\itemsep}{0pt}
\item At each stage $k$, the components $I_{k,1},\ldots,I_{k,2^k-1}$ occur left to right and carry heights $\frac{1}{2^k},\ldots,\frac{2^k-1}{2^k}$.\\
Thus the prescription on $U$ is nondecreasing, and the supremum extension on $P$ preserves order.
\item If $x_0\in U$, then $x_0\in J$ for some component $J$ of $U$ and $\varphi|_J$ is constant. If $x_0\in P$, then for each $k$ there are $a_k,b_k\in2^{-k}\mathbb{Z}\cap[0,1]$ with $a_k\leq\varphi(x_0)\leq b_k$, $b_k-a_k\leq2^{-k}$, and $\varphi(y)\in[a_k,b_k]$ for all $y$ near $x_0$.\\
Thus $\limsup_{y\to x_0}|\varphi(y)-\varphi(x_0)|\leq2^{-k}$ for every $k$, hence $\varphi$ is continuous at $x_0$.
\item The definition gives $\varphi(0)=0$. For each $k$, the rightmost component of $U_k$ has height $1-2^{-k}$, so $\varphi(1)=1$.\\
By (ii) and the intermediate value theorem, $\varphi([0,1])=[0,1]$.
\item For each component $J$ of $U$, $\varphi|_J$ is constant; hence $\varphi'(x)=0$ for every $x\in U$.
\item Proposition \ref{prop:finite-countable-cantor-null-sets} gives $m(P)=0$; since $[0,1]=P\sqcup U$, $m(U)=m([0,1])-m(P)=1$.
\end{enumerate}
\hlt{Structural Commentary}: Plateaus on $U$ force variation onto $P$.\\
\hlt{Key Insight}: At stage $k$, neighboring plateau heights differ by $2^{-k}$.
\end{proof}
\end{proposition}

\begin{theorem} \label{thm:lebesgue-measurable-non-borel-set-exists} \hlt{Lebesgue Measurable Non-Borel Sets Exist}\\
There exist $A\subseteq[0,1]$ and a homeomorphism $\psi:[0,1]\to[0,2]$ such that $A\in\mathcal{L}^1\setminus\mathcal{B}(\R)$ and $\psi(A)\notin\mathcal{L}^1$.
\begin{proof} Let $P$ be the Cantor set, $U\defeq[0,1]\setminus P$, $\varphi$ the Cantor--Lebesgue singular function, and $\psi\defeq\operatorname{id}_{[0,1]}+\varphi$.\\
Proposition \ref{prop:cantor-singular-function-properties} gives $\varphi$ continuous and nondecreasing; hence $\psi$ is continuous and strictly increasing.\\
Since $\psi(0)=0$ and $\psi(1)=2$, $\psi:[0,1]\to[0,2]$ is a homeomorphism.\\
Let $(J_k)_{k\geq1}$ be the components of $U$. Since $\varphi|_{J_k}$ is constant, $\psi(J_k)$ is a translate of $J_k$ and $m(\psi(J_k))=m(J_k)$.\\
Thus $m(\psi(U))=\sum\limits_{k\geq1}m(\psi(J_k))=\sum\limits_{k\geq1}m(J_k)=m(U)=1$.\\
Since $\psi(P)$ is compact and $[0,2]=\psi(P)\sqcup\psi(U)$, $m(\psi(P))=1$.\\
Choose $W\subseteq\psi(P)$ with $W\notin\mathcal{L}^1$ by Theorem \ref{thm:positive-outer-measure-contains-nonmeasurable-subset}, and set $A\defeq\psi^{-1}(W)$. Then $A\subseteq P$ and $m(P)=0$, so $A\in\mathcal{L}^1$. If $A\in\mathcal{B}(\R)$, then $W=\psi(A)\in\mathcal{B}([0,2])\subseteq\mathcal{L}^1$, contradiction. Hence $A\in\mathcal{L}^1\setminus\mathcal{B}(\R)$ and $\psi(A)=W\notin\mathcal{L}^1$.\\
\hlt{Structural Commentary}: Completion by null sets creates Lebesgue sets invisible to the Borel hierarchy.\\
\hlt{Key Insight}: A null Cantor set can be homeomorphically stretched to positive measure.
\end{proof}
\end{theorem}
